% Insert in the Methodology chapter after the prompt-condition overview.
% Requires: booktabs, tabularx, array, ragged2e.

Prior literature supplied the conceptual and methodological building blocks for
this study, but it did not determine the exact wording of the prompts. The
specific prompt formulations, component combinations, and revisions were
implemented by the author. Table~\ref{tab:prompt-revision-logic} distinguishes
the methodological basis of each revision from the contribution made in the
present study.

\begin{table}[ht]
\centering
\small
\renewcommand{\arraystretch}{1.2}
\caption{Methodological basis and author-designed prompt revisions}
\label{tab:prompt-revision-logic}
\begin{tabularx}{\textwidth}{@{}p{0.13\textwidth}X X X@{}}
\toprule
Revision & Problem or motivation & What the literature supports & Author-designed change \\
\midrule
S2 $\rightarrow$ S2$_r$ &
The original instruction specified how to preserve partisan framing but did not explicitly define the treatment of neutral articles. &
Perspective preservation as a summarisation objective (Liu et al., 2024). &
Added a symmetric neutral branch: partisan framing is preserved when present, while neutral articles receive a neutral summary. \\
\addlinespace
S3a $\rightarrow$ S3a$_r$ &
The original condition used dynamically retrieved neutral sentences as a calibration reference. &
BM25 as a relevance-based retrieval method (Robertson and Zaragoza, 2009); framing as a way of analysing emphasis and perspective (Entman, 1993). &
Replaced retrieved examples with an explicit definition of neutral reporting to reduce dependence on cross-dataset retrieval. \\
\addlinespace
S3b $\rightarrow$ S3b$_r$ &
The original prompt presented retrieved sentences as examples of how the same topic is reported across the political spectrum. &
Left--centre--right political ideology as an article-level classification task (Baly et al., 2020); framing theory (Entman, 1993). &
Replaced retrieved spectrum examples with explicit left-, centre-, and right-leaning framing definitions. \\
\addlinespace
S3c $\rightarrow$ S3c$_r$ &
The original condition combined spectrum examples with a faithfulness instruction but did not explicitly specify the neutral case. &
Perspective preservation (Liu et al., 2024) and political-ideology categories (Baly et al., 2020). &
Combined definition-based spectrum guidance with the symmetric neutral branch used in S2$_r$. \\
\addlinespace
S4 $\rightarrow$ S4$_r$ &
The original prompt asked for framing identification but required summary-only output, making the classification step difficult to audit. &
Framing and article-level ideology categories (Entman, 1993; Baly et al., 2020). &
Exposed the framing label in a structured response containing a ``Framing'' field and a ``Summary'' field. \\
\addlinespace
S5 $\rightarrow$ S5$_r$ &
The original two-step prompt did not constrain the stance judgement to a calibrated scale or an explicit textual cue. &
Step-by-step prompting and intermediate reasoning (Wei et al., 2022). &
Added a $-2$ to $+2$ stance scale, required one supporting phrase, and specified a neutral branch for rating 0. \\
\addlinespace
S5$_{adv}$ &
The revised S5 still left the relationship between evidence, stance judgement, and summary implicit. &
Stepwise decomposition (Wei et al., 2022) and rationale-oriented evidence extraction. &
Introduced an evidence-first sequence: extract source phrases, rate the political lean based only on those phrases, and generate the summary. \\
\addlinespace
S6 $\rightarrow$ S6$_r$ &
The original instruction to identify and preserve leaning did not provide a neutral option and could encourage amplification. &
Preservation of the source perspective as a summarisation objective (Liu et al., 2024). &
Distinguished preservation from amplification and added an explicit neutral branch. \\
\bottomrule
\end{tabularx}
\end{table}

The revised prompts should therefore be interpreted as author-designed,
theory-informed interventions rather than direct replications of the cited
studies. Their effectiveness is evaluated empirically in the second-stage
experiment.

\subsection*{Statistical analysis}

Revision effects were tested using paired comparisons on the same 150 articles:
each revised condition was compared with its corresponding original condition
(S2$_r$ vs.~S2, S3a$_r$ vs.~S3a, S3b$_r$ vs.~S3b, S3c$_r$ vs.~S3c,
S4$_r$ vs.~S4, S5$_r$ vs.~S5, S5$_{adv}$ vs.~S5, and S6$_r$ vs.~S6).
Suppression and Injection were analysed with two-sided exact McNemar tests,
and BERTScore with paired two-sided Wilcoxon signed-rank tests.

Because each outcome involved eight simultaneous prompt comparisons, reporting
the raw $p$-values would inflate the probability of at least one false
positive within that outcome family. We therefore applied Holm's step-down
procedure separately to the eight tests for Suppression, the eight tests for
Injection, and the eight tests for BERTScore. Holm correction controls the
family-wise error rate while being less conservative than a one-step
Bonferroni correction (Holm, 1979). We report the resulting adjusted
$p$-values and use 0.05 as the significance threshold.

\noindent\textit{Reference:} Holm, S. (1979). A simple sequentially rejective
multiple test procedure. \textit{Scandinavian Journal of Statistics}, 6(2),
65--70. \url{https://doi.org/10.2307/4615733}.
